% !TEX TS-program = pdflatex
% !TEX encoding = UTF-8 Unicode
%
% Compiled Stateflow: A Typed IR for Out-of-Core Multi-GPU KGE Training
% VLDB submission draft
%
% Class: swap to the official PVLDB class (\documentclass{vldb}) once the
% template ships for the target round.  Using acmart[sigconf] as a
% compatible stand-in while drafting.
\documentclass[sigconf, nonacm]{acmart}

\usepackage{booktabs}
\usepackage{amsmath}
\usepackage{amssymb}
\usepackage{algorithm}
\usepackage{algpseudocode}
\usepackage{listings}
\usepackage{xcolor}
\usepackage{url}

% Lightweight fallback for unit formatting in the draft build.
\providecommand{\SI}[2]{#1\,#2}

% PVLDB copyright placeholder -- replaced by official block at camera-ready.
\settopmatter{printacmref=false}
\renewcommand\footnotetextcopyrightpermission[1]{}
\pagestyle{plain}

\lstset{
  basicstyle=\ttfamily\small,
  keywordstyle=\color{blue}\bfseries,
  commentstyle=\color{gray}\itshape,
  showstringspaces=false,
  columns=fullflexible,
  breaklines=true,
}

\begin{document}

\title{Compiled Stateflow: A Typed IR for Out-of-Core Multi-GPU KGE Training}

\author{Anonymous Authors}
\affiliation{%
  \institution{Anonymous Institution}
  \country{}
}

\begin{abstract}
Knowledge graph embedding (KGE) training at billion-edge scale forces
embeddings and optimizer state out of GPU memory: on a \SI{24}{GB}
GPU, a 41.6M-node graph at 400 dimensions requires \SI{66.6}{GB} of
embeddings and an equivalent amount of Adagrad state. State-of-the-art
out-of-core trainers address single-GPU residency through partitioned
schedules (e.g., the COVER schedule used in GE\textsuperscript{2}), but
extend to multi-GPU settings by treating each device as an independent
COVER instance. This leaves inter-GPU handoffs implicit, hides their
cost from the scheduler, and forecloses peer-to-peer transfers whose
cost model differs fundamentally from host reloads.
\par\vspace{0.5em}
We introduce \emph{Compiled Stateflow}, a typed intermediate
representation (IR) for out-of-core partition schedules. Stateflow makes
every byte of partition residency, host admit, and cross-device handoff
a first-class object with an explicit cost. A lane-matching compiler
chooses per-GPU assignments that minimize host I/O under a unified cost
model, exposing cross-lane peer opportunities as typed
\texttt{PeerHandoff} descriptors executable as \texttt{cudaMemcpyPeerAsync}.
On Twitter (41.6M nodes, 1.46B edges) with 16 partitions, our compiler
reduces compile-time selection cost by 14\% and eliminates all
\texttt{FULL\_RELOAD} handoffs, exposing \SI{37.49}{GB} of per-epoch
peer-transfer opportunity. On a two-GPU RTX A5000 PCIe system, the
runtime executes the embedding half of that opportunity
(\SI{18.74}{GB}/epoch) with zero host fallback and zero descriptor
mismatches. Epoch time remains effectively flat relative to host
reloads, showing that the IR and runtime are correct even when this
interconnect does not turn peer copies into a wall-clock win. We
present the IR, cost model, lane-matching algorithm, validator
invariants, and measured runtime behavior.
\end{abstract}

\keywords{knowledge graph embeddings, out-of-core training, multi-GPU
scheduling, intermediate representations, peer-to-peer transfers}

\maketitle

%-----------------------------------------------------------------------
\section{Introduction}\label{sec:intro}
% Section 1 -- Introduction
%
% Shape:
%   1.1 Motivating gap: KGE embeddings no longer fit in GPU memory,
%       and the fix for a single GPU does not generalize to many.
%   1.2 Why multi-GPU COVER-replication is the wrong baseline.
%   1.3 Thesis: the scheduler needs a typed, costed IR.
%   1.4 Contributions + results preview.

Knowledge graph embedding (KGE) training pushes two tensors through
the GPU on every step: a block of embedding rows and a matching block
of optimizer state. Both scale linearly with the number of nodes and
with the embedding dimension $d$, and both must be writable on the
device that touches them. At today's dataset sizes that product no
longer fits. A Twitter-scale graph with 41.6\,M nodes at $d\!=\!400$ in
fp32 requires \SI{66.6}{GB} for the embedding matrix and an equal
\SI{66.6}{GB} for Adagrad's accumulator -- a combined \SI{133}{GB}
against a \SI{24}{GB} device-memory budget. Billion-edge KGE
has been an \emph{out-of-core} workload for several years now, and the
community has responded with partitioned trainers that stream
residency in and out of the GPU on a schedule.

\paragraph{The single-GPU story has stabilized.}
Recent work on partitioned KGE -- GE\textsuperscript{2}~\cite{ge2_report},
Marius~\cite{marius}, DGL-KE~\cite{dgl_ke}, PyTorch-BigGraph~\cite{pytorch_biggraph} --
has converged on a common recipe: partition the node set into $p$
blocks, schedule microsteps over the resulting $p^2$ edge buckets so
that each block is resident when its edges are trained, and keep the
blocks on the GPU long enough to amortize loading. GE\textsuperscript{2}'s
COVER algorithm formalizes this as a Partition Scheduling Problem and
reduces the schedule to a Resolvable Balanced Incomplete Block Design
(RBIBD) with parameters $(p, q, 1)$. The result is optimal in admit
count for single-GPU execution and matches the communication
lower-bound $\tfrac{4GNd}{3}$ to within a small constant factor.

\paragraph{The multi-GPU extension is not a schedule.}
The natural way to use two GPUs with a COVER-based trainer is to run
two independent COVER instances, each on a disjoint subset of rounds.
This is what every current partitioned KGE system actually does.
It scales throughput approximately linearly in the common case, but
it makes two decisions silently:

\begin{enumerate}
\item \emph{Inter-GPU handoffs are charged as host reloads.} If lane
0 held partition $\pi$ in round $r$ and lane 1 needs $\pi$ in round
$r\!+\!1$, the system evicts $\pi$ from lane 0 and re-admits it to
lane 1 from the host -- even though $\pi$'s most recent authoritative
copy lives one PCIe hop away on a peer device. The peer path is
invisible to the scheduler.
\item \emph{Lane assignment is a side effect of the partition
numbering.} Because each lane runs its own COVER, the lane a
partition lands on in round $r$ is determined by the order COVER
enumerates microsteps, not by any attempt to keep a partition on the
same lane across rounds. A favorable assignment happens by accident;
an unfavorable one happens as often.
\end{enumerate}

Both are scheduling decisions, but the multi-GPU system has nowhere to
express them. GE\textsuperscript{2}'s scheduler is a COVER solver; its
output is a list of per-device buckets and a host reload pattern. It
does not carry a notion of a \emph{cross-lane handoff}, so it has no
way to reason about one, optimize against one, or hand one to a
runtime that could execute it peer-to-peer.

\paragraph{Thesis.}
We argue that the multi-GPU gap is a representation problem before it
is a performance problem. The scheduler is not missing an algorithm;
it is missing a type. If a compiler can see that a partition is
resident on lane 0 at round $r$ and needed on lane 1 at round $r\!+\!1$,
it can choose whether to serve that admit from host memory or from
a peer device, weigh the two against each other under an explicit
cost model, and commit to one. Today's systems cannot even formulate
the question, because their output format has no field for the
answer.

Our contribution is \emph{Compiled Stateflow}, a typed intermediate
representation (IR) for out-of-core partition schedules, and a
lane-matching compiler that produces plans in that IR. Stateflow
makes every byte of partition residency, host admit, and cross-device
handoff a first-class value with an explicit cost. A plan in
Stateflow is a data structure -- produced, validated, and scored
offline -- that carries four distinct handoff modes
(\texttt{FULL\_RELOAD}, \texttt{ROTATING\_OVERWRITE},
\texttt{DELAYED\_KEEP\_ALIVE}, \texttt{PEER\_RELAY}) and a
decomposable cost function over admits, bucket balance, boundary
crossings, and lane imbalance. The runtime consumes
\texttt{PeerHandoffDescriptor}s projected from the plan and can
execute them as \texttt{cudaMemcpyPeerAsync} without any scheduling
logic of its own.

\paragraph{Results preview.}
On Twitter (41.6\,M nodes, 1.46\,B edges), a 16-partition configuration
on two \SI{24}{GB} RTX A5000s, our lane-matching compiler produces
plans that:
\begin{itemize}
\item lower the host-denominated selection cost from
$1.499 \!\times\! 10^{11}$ to $1.291 \!\times\! 10^{11}$ -- a
\textbf{14\% reduction} in compile-time cost,
\item cut host admits from 72 to 62 (14\%),
\item \textbf{eliminate every \texttt{FULL\_RELOAD} handoff}
(10 $\rightarrow$ 0), converting them into 18
\texttt{ROTATING\_OVERWRITE} and 18 \texttt{PEER\_RELAY} descriptors,
\item expose \textbf{\SI{37.49}{GB} of planned device-to-device traffic}
for the runtime, versus \SI{16.66}{GB} for
\textsc{Disjoint\_Rounds},
\item and, in the final embeddings-scope runtime, execute
\textbf{all 18 planned peer relays} for \textsc{Lane\_Matched}
(\SI{18.74}{GB}, zero fallback, zero descriptor mismatch), while
leaving epoch time effectively flat
(\SI{209.450}{s} vs.\ \SI{208.256}{s} for the host-reload baseline).
\end{itemize}

The compile-time wins are attributable to the compiler; the runtime
result shows that the projected descriptors are now executable and
measurable end-to-end. On this PCIe machine, Phase 4b is
correctness-positive but not yet wall-clock-positive.

\paragraph{Contributions.}
\begin{enumerate}
\item \textbf{Compiled Stateflow}, a typed IR for out-of-core
partition scheduling that separates intra-lane from cross-lane
handoffs, makes every transfer a typed object, and supports a
decomposable cost model (Section~\ref{sec:ir}).
\item \textbf{A lane-matching compiler} for multi-GPU KGE training
that emits Stateflow plans with cross-lane \texttt{PEER\_RELAY}
descriptors, under a host-denominated cost model that does not
credit itself by default for peer savings that are reported but not
optimized
(Section~\ref{sec:compiler}).
\item \textbf{A five-invariant validator} (V1--V5) that rejects
structurally inconsistent or unexecutable plans before they reach
the runtime (Section~\ref{sec:ir:validator}).
\item \textbf{A peer-aware runtime path} built around
\texttt{cudaMemcpyPeerAsync}, gated behind a kill switch and a
fallback to host reloads when peer access is unavailable
(Section~\ref{sec:runtime}).
\item \textbf{An evaluation} on Twitter (41.6\,M, 1.46\,B) showing
the compiler's compile-time gains and the runtime's exact zero-fallback
execution of planned peer relays; on our PCIe A5000 setup, the
peer-enabled path is effectively tied with the host baseline rather
than faster (Section~\ref{sec:eval}).
\end{enumerate}

The rest of the paper is organized as follows.
Section~\ref{sec:background} reviews KGE training, out-of-core
partitioning, and the COVER baseline.
Section~\ref{sec:ir} introduces the Stateflow IR, its types, cost
model, and validator.
Section~\ref{sec:compiler} presents the lane-matching compiler.
Section~\ref{sec:runtime} describes the peer-aware runtime.
Section~\ref{sec:eval} evaluates both on Twitter.
Section~\ref{sec:related} relates Stateflow to other IRs and other
out-of-core trainers, and Section~\ref{sec:conclusion} concludes.


%-----------------------------------------------------------------------
\section{Background}\label{sec:background}
% Section 2 -- Background  [DRAFT STUB]
%
% Intended beats:
%   - KGE training model: scorer, negative sampling, Adagrad state
%     doubling the resident footprint.
%   - Out-of-core partitioning: why node embeddings are the bottleneck,
%     not edges.
%   - GE^2 / COVER: Partition Scheduling Problem as RBIBD(p, q=4, 1);
%     Theorem 1/Corollary 1; communication lower bound 4*G*N*d / 3.
%   - Multi-GPU in GE^2: independent COVER instances per device;
%     no cross-device scheduling.
%   - Peer interconnects: PCIe 4.0 vs NVLink; cudaMemcpyPeerAsync
%     semantics; why the host path dominates without explicit peer
%     scheduling.

\emph{[Background draft pending.]}


%-----------------------------------------------------------------------
\section{Compiled Stateflow IR}\label{sec:ir}
% Section 3 -- Compiled Stateflow IR
%
% Goal: introduce the typed IR.  Keep it self-contained -- a reader who
% skipped the background should still be able to see what a plan is,
% what it costs, and how we reject bad plans.

\subsection{Design Goals}\label{sec:ir:goals}

Out-of-core KGE training is a scheduling problem whose decisions --
which partition to evict, which to admit, from where to source the
admit -- happen at runtime but are determined by choices made long
before any GPU kernel runs. We designed Compiled Stateflow around four
goals that fall out of that observation.

\paragraph{G1: Every byte is typed.} Host admits, device-to-device peer
copies, and kept-alive resident pages differ by two or three orders of
magnitude in latency and by an order of magnitude in bandwidth. A
scheduler that cannot distinguish them cannot reason about them. Every
transfer in Stateflow carries a mode tag and a byte count.

\paragraph{G2: Plans are values, not side effects.} A Stateflow plan is
a pure data structure. It is produced by the compiler, handed to a
validator, and only then consumed by the runtime. This separation lets
us enumerate candidates, cost them, apply invariants, and diff plans
across configurations without running the trainer.

\paragraph{G3: Cost is explicit and decomposable.} The scheduler's
objective function is not a scalar black box. A plan's cost decomposes
into admit cost, bucket imbalance, boundary crossings, and lane
imbalance, each attributable to a concrete portion of the plan. This
lets the compiler explain \emph{why} one plan beat another, which in
turn lets us write targeted transformations.

\paragraph{G4: The IR survives Phase changes.} We stage the system in
phases -- single-GPU scheduling, multi-GPU compilation, peer runtime,
bandwidth-aware staggering -- and the IR is the one artifact that
remains stable across them. Phase 3 adds descriptors the runtime does
not yet execute; the final Phase 4b runtime now executes them without
changing the compiler. Phase 5's staggering will introduce new cost
terms without changing the type hierarchy. A type that absorbs
additions without breaking consumers is the whole reason to build an
IR at all.

\subsection{Type Hierarchy}\label{sec:ir:types}

A plan is a nested structure (Figure~\ref{fig:ir:types}). At the top,
\texttt{StateflowPlan} holds one \texttt{LanePlan} per active GPU plus
a top-level list of cross-lane handoffs. Each \texttt{LanePlan} is an
ordered sequence of \texttt{Microstate}s -- the per-lane steps between
round boundaries -- connected by intra-lane \texttt{Handoff}s. A
\texttt{Microstate} names the partitions resident on the device at that
step, the fragments they occupy, and the resident objects (embeddings,
optimizer state) that back them.

\begin{figure}[t]
\centering
\begin{lstlisting}[basicstyle=\ttfamily\footnotesize]
StateflowPlan
  family        : {CUSTOM, HYBRID_COVER, MULTI_GPU}
  variant       : {CANONICAL, GPU_AWARE, REVERSED, ...
                   DISJOINT_ROUNDS, LANE_MATCHED}
  lanes         : [LanePlan]
  cross_lane_handoffs : [HandoffPlan]    # top-level only
  cost          : PlanCostBreakdown
  total_peer_handoff_bytes : int64

LanePlan
  lane_id       : int
  microstates   : [MicrostatePlan]       # length N
  handoffs      : [HandoffPlan]          # length N-1, intra-lane

MicrostatePlan
  round_idx     : int
  resident_partitions : [int]
  fragments     : [FragmentPlan]
  resident_objects    : [ResidentObjectPlan]

HandoffPlan
  mode          : {FULL_RELOAD, ROTATING_OVERWRITE,
                   DELAYED_KEEP_ALIVE, PEER_RELAY}
  partition_id  : int
  src_lane_id   : int     # -1 if host-sourced
  dst_lane_id   : int
  bytes         : int64
  peer_bytes    : int64   # nonzero iff PEER_RELAY

PeerHandoffDescriptor                    # runtime target
  src_lane, dst_lane, src_slot, dst_slot : int
  round_idx, partition_id                : int
  bytes                                  : int64
\end{lstlisting}
\caption{Core Stateflow types. The IR separates intra-lane handoffs
(attached to a \texttt{LanePlan}) from cross-lane handoffs (attached to
the top-level plan). The runtime consumes \texttt{PeerHandoffDescriptor}s
projected from the plan.}
\label{fig:ir:types}
\end{figure}

Two design choices in Figure~\ref{fig:ir:types} are worth calling out.

First, \textbf{cross-lane handoffs live at the top level}, not inside
\texttt{LanePlan.handoffs}. A \texttt{LanePlan} with $N$ microstates
must have exactly $N-1$ intra-lane handoffs; this invariant is checked
by the validator and lets downstream code zip lanes with their own
transitions. Mixing cross-lane handoffs into that list would break the
invariant or require per-entry tagging at every call site. We instead
expose them through \texttt{StateflowPlan.cross\_lane\_handoffs}, a
flat list the runtime indexes by \texttt{round\_idx}.

Second, \texttt{HandoffPlan} carries both \texttt{bytes} and
\texttt{peer\_bytes}. The former is the host-fallback cost (what a
host admit would transfer); the latter is nonzero only for
\texttt{PEER\_RELAY} and captures the device-to-device byte count the
runtime would issue via \texttt{cudaMemcpyPeerAsync}. Keeping them
separate lets the default cost model be host-denominated (safe) while
still reporting peer-executable volume for Phase 4 headroom analysis.

\subsection{Handoff Modes}\label{sec:ir:modes}

Partition residency changes at round boundaries. We identified four
distinct ways a partition can become resident on lane $\ell$ at round
$r{+}1$; each is a separate handoff mode because each has a different
cost and different runtime mechanism.

\begin{description}
\item[\texttt{FULL\_RELOAD}] The partition is fetched from host memory
as if this were a cold admit. Cost: full host-bandwidth transfer of
embeddings and optimizer state. Occurs when no lane held the partition
in round $r$, or when the holder did not share an NVLink/PCIe peer
path to $\ell$.
\item[\texttt{ROTATING\_OVERWRITE}] The partition was resident on $\ell$
in round $r$ and its slot is reused in round $r{+}1$. No transfer; the
cost is zero. Named for the rotating-window layouts that naturally
produce this pattern.
\item[\texttt{DELAYED\_KEEP\_ALIVE}] The partition was resident on $\ell$
in round $r$, is not needed in round $r{+}1$, but is scheduled for
round $r{+}2$. The runtime keeps it pinned to avoid reloading it one
round later. Cost: device memory occupancy for one extra round; no
transfer.
\item[\texttt{PEER\_RELAY}] The partition is resident on lane $\ell'$
in round $r$, not on $\ell$, but needed on $\ell$ in round $r{+}1$. A
device-to-device copy from $\ell'$ to $\ell$ delivers it. Cost:
\texttt{peer\_bytes} at peer bandwidth (substantially below host
bandwidth on modern interconnects). Requires
\texttt{cudaDeviceCanAccessPeer} to be true.
\end{description}

\texttt{PEER\_RELAY} is the mode that makes the multi-GPU case
interesting. The compiler emits it whenever lane matching produces a
cross-lane opportunity; the runtime executes it as
\texttt{cudaMemcpyPeerAsync}. The same descriptor supported both
bring-up regimes: Phase 3 logged it while still executing a host
reload, and the final Phase 4b runtime consumes it directly without
changing the compiler or validator.

\subsection{Cost Model}\label{sec:ir:cost}

Candidate plans are ranked by a weighted sum of four terms:
\begin{align}
\mathrm{cost}(P) =\ & \alpha \cdot \mathrm{admit}(P)
                    + \beta  \cdot \mathrm{bucket}(P) \nonumber \\
                  & + \gamma \cdot \mathrm{boundary}(P)
                    + \delta \cdot \mathrm{imbalance}(P).
\label{eq:cost}
\end{align}

\paragraph{Admit cost.} $\mathrm{admit}(P) = \sum_h b_h$, summed over
handoffs $h$ that cause a host fetch. Measured in bytes; this is the
host-bandwidth-denominated cost of populating microstates whose
partitions were not resident in the preceding round. In default
configuration, \texttt{PEER\_RELAY} is counted as a \emph{host} admit
(its bytes would be loaded from host if the peer path were
unavailable). This is a deliberately conservative choice: until the
selector is made peer-aware, the compiler must not give itself credit
for savings it is not yet optimizing for.

\paragraph{Bucket cost.} $\mathrm{bucket}(P)$ penalizes uneven edge
bucket sizes across microstates. Small buckets underutilize the GPU;
enormous buckets over-fill a single superstep. We weight this term
with a small $\beta \!\approx\! 10^{-7}$ so it acts as a tie-breaker
among plans with similar admit cost.

\paragraph{Boundary cost.} $\mathrm{boundary}(P)$ counts superstate
boundaries -- transitions that force global synchronization across
lanes. In the multi-GPU setting these are NCCL all-reduce points;
each one is a latency cost plus a bandwidth tax on embedding
gradients.

\paragraph{Imbalance cost.} $\mathrm{imbalance}(P) = \max_\ell W_\ell
- \min_\ell W_\ell$, where $W_\ell$ is lane $\ell$'s per-round work
(typically summed bucket bytes). A plan where one lane finishes
substantially before another wastes the gap; this term pushes the
compiler toward lane assignments where the work lands evenly.

\paragraph{Reported metrics, not optimized.} Two quantities are
recorded on every \texttt{StateflowPlan} but not summed into
Equation~\ref{eq:cost}: \texttt{total\_peer\_handoff\_bytes} and
\texttt{total\_cross\_lane\_handoffs}. These characterize the plan's
Phase 4 opportunity (how much host I/O a peer runtime could displace)
without letting the selector claim credit for it during compilation.
Even now that the runtime executes peer copies, the default selector
keeps these as reported metrics rather than optimized terms. A future
peer-aware cost model can sum them into Equation~\ref{eq:cost}
explicitly -- a clean extension point that does not disturb Phases 1--4.

\subsection{Validator Invariants}\label{sec:ir:validator}

A plan is only useful if it is executable. The validator rejects plans
that violate any of the following invariants; no plan leaves the
compiler without passing.

\begin{description}
\item[V1: Structural consistency.] Each lane has
$|\text{handoffs}| = |\text{microstates}| - 1$. Cross-lane handoffs
appear only in \texttt{StateflowPlan.cross\_lane\_handoffs}, never in
\texttt{LanePlan.handoffs}.
\item[V2: Coverage.] For every microstate boundary $(r, r{+}1)$ on
every lane $\ell$, the resident set at $r{+}1$ equals
$(\text{kept}_{\ell}) \cup (\text{host-admitted}) \cup
 (\text{peer-received})$. No partition appears on a lane without a
sourcing handoff.
\item[V3: Peer admit sourcing.] A partition delivered by
\texttt{PEER\_RELAY} must be resident on the source lane in round $r$
and must not be resident on the destination lane in round $r$. Both
ends are checked.
\item[V4: Superstate disjointness.] Within a round, partitions
assigned to the same superstate must be disjoint across lanes.
Violating this would double-count gradient contributions during the
all-reduce.
\item[V5: Handoff consistency.] Every cross-lane handoff with mode
\texttt{PEER\_RELAY} has matching source and destination lanes in the
microstate's resident sets. The byte count agrees with the embedding
layout (Section~\ref{sec:ir:layout}).
\end{description}

Invariants V1--V5 are machine-checked and cheap -- the validator
passes in microseconds on 16-partition plans. They serve two
purposes: they catch compiler bugs at their source (a corrupt plan
cannot reach the runtime), and they document the contract between the
IR and every downstream consumer.

\subsection{Embedding Layout}\label{sec:ir:layout}

Byte counts are not inferable from partition IDs alone: they depend on
embedding dimension, dtype, and whether the trainer uses optimizer
state that doubles the resident footprint. \texttt{PlanEmbeddingLayout}
captures this:

\begin{lstlisting}[basicstyle=\ttfamily\footnotesize]
PlanEmbeddingLayout
  embedding_dim           : int    # e.g. 400
  dtype_size              : int    # bytes per element
  optimizer_state_multiplier : int # 1 or 2
\end{lstlisting}

\noindent
A partition of $n$ rows contributes
$n \cdot \mathtt{embedding\_dim} \cdot \mathtt{dtype\_size}
     \cdot \mathtt{optimizer\_state\_multiplier}$ bytes to any handoff
that transfers it. The layout is populated once at plan construction
from the graph storage metadata, then embedded in every plan so
downstream tooling (cost reporters, visualizers) can produce concrete
byte figures without re-querying the model.

\subsection{A Concrete Example}\label{sec:ir:example}

Consider Twitter (41.6M nodes, $d\!=\!400$, fp32, Adagrad), 16
partitions, buffer capacity 4, two GPUs. Two compiled plans:

\begin{table}[h]
\centering
\small
\begin{tabular}{lrr}
\toprule
 & \textsc{Disjoint\_Rounds} & \textsc{Lane\_Matched} \\
\midrule
selection cost        & $1.499 \times 10^{11}$ & $1.291 \times 10^{11}$ \\
host admits           & 72                     & 62                     \\
cross-lane handoffs   & 8                      & 18                     \\
\texttt{PEER\_RELAY} bytes & 16.66 GB          & 37.49 GB               \\
\texttt{FULL\_RELOAD} count & 10               & 0                      \\
\texttt{ROTATING\_OVERWRITE} & 8               & 18                     \\
boundaries            & 18                     & 18                     \\
\bottomrule
\end{tabular}
\caption{Two compiled Stateflow plans on Twitter 16-partition, 2-GPU.
\textsc{Lane\_Matched} lowers host-denominated cost by 14\% and
eliminates all \texttt{FULL\_RELOAD} events and exposes 37.49 GB of
peer-relay opportunity. Section~\ref{sec:eval} shows the runtime
executing the embeddings half of that traffic with zero fallback.}
\label{tab:ir:example}
\end{table}

Both plans satisfy V1--V5 and are produced by the same compiler under
different variant flags. \textsc{Lane\_Matched} is the output of the
compiler described in Section~\ref{sec:compiler}; \textsc{Disjoint\_Rounds}
is a baseline that ignores cross-lane opportunities. The entire
difference is visible in the IR -- no runtime is needed to see that
\textsc{Lane\_Matched} has eliminated the full reloads. The runtime
results of Section~\ref{sec:eval} then show that these IR-level
differences survive projection and swap-time execution exactly.


%-----------------------------------------------------------------------
\section{Lane-Matching Compiler}\label{sec:compiler}
% Section 4 -- Lane-Matching Compiler

\subsection{Compiler Overview}\label{sec:compiler:overview}

The lane-matching compiler turns a single-GPU COVER/Hybrid-Cover
schedule into a multi-lane Stateflow plan. Its input is a list of
\emph{buffer states} -- each state is a $q$-subset of partitions that
would occupy one device for one microstep on a single GPU -- together
with the per-state edge buckets that cover the resulting edges. Its
output is a permutation of those states arranged into rounds of
$k$ lanes, along with a populated \texttt{StateflowPlan} including
cross-lane \texttt{PEER\_RELAY} descriptors.

The compiler runs in three passes (Figure~\ref{fig:compiler:pipeline}).
First, \emph{lane assignment} chooses, for each round, which buffer
states run on which lane so as to maximize residency overlap with the
previous round. Second, \emph{plan construction} materializes the
permutation into a \texttt{StateflowPlan} with per-lane microstates and
intra-lane handoffs. Third, \emph{cross-lane lifting} walks the plan
and emits a \texttt{PeerHandoffDescriptor} whenever a microstate
admits a partition that a peer lane held in the previous round.

\begin{figure}[t]
\centering
\begin{lstlisting}[basicstyle=\ttfamily\footnotesize]
buffer_states, edge_buckets  (single-GPU COVER output)
        |
        v
[Pass 1] assign_lanes_per_round  -> permutation: [s0,s1,s2,s3,...]
        |                            (round r uses permutation[r*k..r*k+k])
        v
[Pass 2] build_multi_gpu_stateflow_plan_from_permutation
        |   ->  plan with per-lane microstates + intra-lane handoffs
        v
[Pass 3] populate_cross_lane_handoffs
            ->  plan.cross_lane_handoffs += PEER_RELAY descriptors
        |
        v
[Pass 4] score_stateflow_plan    (selection_cost, planned_peer_bytes)
        |
        v
[Pass 5] validator (V1--V5)      (reject or accept)
\end{lstlisting}
\caption{Lane-matching compiler pipeline. Passes 1--3 produce the
plan; Pass 4 scores it; Pass 5 rejects it on invariant violation.}
\label{fig:compiler:pipeline}
\end{figure}

The rest of this section fills in the two non-trivial passes. Passes
2, 4, and 5 are mechanical: Pass~2 re-packs tensors into microstates
(Section~\ref{sec:ir:types}); Pass~4 is Equation~\ref{eq:cost}; Pass~5
is the validator of Section~\ref{sec:ir:validator}.

\subsection{Pass 1: Per-Round Lane Assignment}\label{sec:compiler:assign}

At each round $r$, the compiler must choose $k$ buffer states from
those not yet placed, and it must order them across lanes
$\ell \in \{0, \dots, k-1\}$. Two states may be placed in the same
round only if they are \emph{disjoint} -- the intersection of their
partition sets is empty -- otherwise they would double-load the same
partition on two devices and break superstate disjointness (invariant
V4). Compatibility is precomputed as a symmetric boolean matrix over
all state pairs.

Given disjointness as a hard constraint, the assignment problem at
round $r$ is:
\[
\arg\min_{(a_0, \dots, a_{k-1})}
  \underbrace{\sum_{\ell=0}^{k-1}
              T(\mathrm{prev}_\ell,\, a_\ell)}_{\text{transition cost}}
  \ +\
  \underbrace{I(a_0, \dots, a_{k-1})}_{\text{imbalance penalty}},
\]
where $a_\ell$ is the state assigned to lane $\ell$,
$\mathrm{prev}_\ell$ is what lane $\ell$ held in round $r\!-\!1$, and
the sum runs over all disjoint $k$-tuples drawn from the remaining
states.

\paragraph{Transition cost.} The per-lane term $T(p, n)$ prices the
handoff from previous state $p$ to next state $n$ on the same lane:
\begin{equation}
T(p, n) = \frac{\mathrm{host\_bytes}(p, n)}{B_{\mathrm{host}}}
         + [\mathrm{overlap}(p, n) = 0] \cdot w_{\mathrm{bnd}}.
\label{eq:transition}
\end{equation}
Here $\mathrm{host\_bytes}(p, n)$ is the byte volume that would need
to be loaded from host if the admit were served by \texttt{FULL\_RELOAD}
(partitions in $n$ but not in $p$, scaled by the embedding layout);
$B_{\mathrm{host}}$ is the host bandwidth from
\texttt{LaneMatchCostConfig}; and $w_{\mathrm{bnd}}$ is an additive
boundary penalty applied whenever residency changes completely between
rounds (the typical case for disjoint-round baselines).

Equation~\ref{eq:transition} is intentionally \emph{host-denominated}.
In Phase 3, peer transfers are planned but not executed, so crediting
a lane assignment with peer-bandwidth savings it will not realize
would make the compiler over-eager. A configuration flag
(\texttt{allow\_peer\_relay}) enables an optional $-\,\mathrm{peer\_bytes}
/B_{\mathrm{peer}}$ credit term; it is disabled by default and guarded
behind \texttt{GEGE\_STATEFLOW\_ALLOW\_PEER\_RELAY=1}. We leave it off
for the headline measurements.

\paragraph{Imbalance penalty.} The per-round term $I$ discourages
assignments where one lane is handed a large bucket and the other a
small one, which would waste the gap between their respective
completions:
\begin{equation}
I(a_0, \dots, a_{k-1}) = w_{\mathrm{imb}}
  \cdot \frac{\sum_\ell (E_{a_\ell} - \bar E)^2}{\bar E},
\label{eq:imbalance}
\end{equation}
where $E_{a_\ell}$ is lane $\ell$'s total edge-bucket volume for the
round and $\bar E$ is the mean across lanes. This is a chi-square-style
normalization: the penalty scales with relative, not absolute,
imbalance. $w_{\mathrm{imb}}$ defaults to $0.25$.

\subsection{Solvers}\label{sec:compiler:solvers}

The per-round problem has $\binom{m}{k} \cdot k!$ candidate
$k$-tuples, where $m$ is the number of unplaced states. Three solvers
are provided; the choice is controlled by
\texttt{GEGE\_STATEFLOW\_LANE\_MATCH\_SOLVER}.

\begin{description}
\item[\texttt{optimal2} (default).] Historical name for the exact
small-$k$ solver. The compiler enumerates all compatible disjoint
$k$-tuples and, for each tuple, solves the lane-assignment subproblem
exactly with subset DP over the $k \times k$ transition-cost matrix.
At $k{=}2$ this degenerates to the original exhaustive pair search;
at $k{=}4$ it remains easily tractable and preserves the same
host-bytes / boundary / imbalance objective as the 2-lane case.

\item[\texttt{greedy}.] Incremental search that places states one at a
time, choosing each to maximize overlap with the corresponding
previous-round lane. Cheap, but it can commit early to an assignment
that is locally good and globally suboptimal. Used as a fallback when
the exact solver is intentionally disabled.

\item[\texttt{hungarian} (alias for \texttt{optimal2}).] Accepted for
backward compatibility with earlier drafts. The current implementation
routes it to the same exact small-$k$ solver.
\end{description}

At two GPUs -- the focus of this paper -- exhaustive enumeration is
both optimal and trivially auditable. The same exact objective now
extends unchanged to four GPUs; the only difference is that the
assignment step is solved with subset DP instead of a hand-written
two-lane special case. \texttt{optimal2} remains the default.

\subsection{Pass 3: Cross-Lane Lifting}\label{sec:compiler:lift}

Once the per-round assignment is baked into a permutation and
materialized by Pass 2, the resulting plan already has correct
per-lane handoffs (each \texttt{LanePlan} has $N$ microstates and $N-1$
intra-lane \texttt{HandoffPlan}s). It does \emph{not} yet know which
admits could have been peer-served; that information is computed in
Pass 3 by walking the plan.

\begin{algorithm}[t]
\caption{Cross-lane lifting (\texttt{populate\_cross\_lane\_handoffs})}
\label{alg:lift}
\begin{algorithmic}[1]
\State $\text{cross\_lane\_handoffs} \gets \emptyset$
\For{each lane $\ell$ and each microstate $m$ on $\ell$ with $m \ne m_0$}
  \State let $m_{-1}$ be the previous microstate on $\ell$
  \For{each partition $\pi \in m.\mathrm{resident\_partitions}$ not in $m_{-1}$}
    \State $\text{best} \gets \bot$
  \For{each other lane $\ell'$}
      \State let $m'$ be $\ell'$'s microstate at superstate $m.\mathrm{superstate\_id}{-}1$
      \If{$m' \ne \bot$ and $\pi \in m'.\mathrm{resident\_partitions}$}
        \State $b \gets \mathrm{partition\_transfer\_bytes}(\pi)$
        \If{$\text{best} = \bot$ \textbf{or} $\ell' < \text{best.lane}$}
          \State $\text{best} \gets (\ell', m', b)$
        \EndIf
      \EndIf
    \EndFor
    \If{$\text{best} \ne \bot$}
      \State emit \texttt{PEER\_RELAY} handoff from $\text{best}.\ell'$ to $\ell$, partition $\pi$
      \State $\text{plan.total\_peer\_handoff\_bytes} \mathrel{+}= \text{best.bytes}$
    \EndIf
  \EndFor
\EndFor
\end{algorithmic}
\end{algorithm}

Algorithm~\ref{alg:lift} is complete: for every admit on every lane,
it either finds a peer source (in which case a \texttt{PeerHandoff}
descriptor is emitted and accounted against
\texttt{total\_peer\_handoff\_bytes}) or it does not (in which case
the admit will be served by host fallback at runtime, already represented
by the intra-lane \texttt{FULL\_RELOAD} in Pass 2). Deduplication by
$(\ell, m, \pi)$ prevents double-counting when a partition is
simultaneously resident on multiple peer lanes.

\paragraph{Source tie-breaking.} Under the current byte model,
\texttt{partition\_transfer\_bytes} depends only on the partition being
transferred, not on which peer lane serves it. When two peer lanes
could each service the same admit, we therefore break ties by smaller
lane id. The rule is deliberately simple and deterministic; a future
bandwidth-aware model could replace it with link-specific source
selection without changing the rest of the pass.

\paragraph{What Pass 3 does \emph{not} do.} It does not modify the
lane assignment chosen by Pass 1. The per-lane permutation is frozen
by the time Pass 3 runs. Co-designing the two passes -- letting the
availability of a peer source pull the per-lane assignment toward
configurations that maximize such sources -- is a natural next step.
The cleaner staging is to run them independently, confirm the per-lane
story is correct, then extend.

\subsection{Variants and Baseline}\label{sec:compiler:variants}

The compiler produces two named multi-GPU variants, both of which
share Passes 2--5:

\begin{description}
\item[\textsc{DISJOINT\_ROUNDS}.] A baseline that skips Pass 1 entirely
and routes to \texttt{getDisjointBufferStatePermutation}: states are
partitioned into disjoint rounds in their COVER-enumeration order,
then assigned to lanes in round order. This is the behavior of running
two independent COVER schedulers on alternating rounds -- the status
quo for multi-GPU COVER-based trainers (Section~\ref{sec:intro}).
\item[\textsc{LANE\_MATCHED}.] The full Pass 1 + Pass 3 pipeline
described above. This is the variant the compiler would choose under
default scoring on Twitter (Section~\ref{sec:eval}).
\end{description}

Both variants satisfy V1--V5; both populate
\texttt{cross\_lane\_handoffs} and \texttt{total\_peer\_handoff\_bytes}.
The difference is entirely in the lane assignment. On Twitter 16p
2-GPU, this difference alone changes the handoff-mode distribution
from $\{\texttt{FULL\_RELOAD}:10,\ \texttt{ROTATING\_OVERWRITE}:8,\
\texttt{PEER\_RELAY}:8\}$ to
$\{\texttt{FULL\_RELOAD}:0,\ \texttt{ROTATING\_OVERWRITE}:18,\
\texttt{PEER\_RELAY}:18\}$, eliminating every cold admit from the plan.

\subsection{Pass 4: Scoring}\label{sec:compiler:score}

Scored plans expose every term of Equation~\ref{eq:cost} on
\texttt{StateflowPlan.cost\_breakdown}:
\texttt{admitted\_partition\_cost},
\texttt{bucket\_edge\_cost},
\texttt{boundary\_cost},
\texttt{lane\_imbalance\_cost}, and the summed
\texttt{selection\_cost}. For multi-GPU plans, the admitted term is
byte-accurate (summed $\mathrm{rows}(\pi) \cdot \mathrm{bytes\_per\_row}$
over admitted objects); for single-GPU plans, it falls back to a
count-based weighted admission load (preserving backward compatibility
with existing single-GPU cost reports).

\texttt{planned\_peer\_bytes} is populated alongside
\texttt{selection\_cost} but \emph{not summed into it}. This is the
same split introduced in Section~\ref{sec:ir:cost}: the compiler
reports the Phase 4 opportunity without letting itself claim credit
for it. A consumer that wants to optimize under a peer-aware cost
model can sum the two terms explicitly; the default pipeline does not.

\subsection{Complexity and Determinism}\label{sec:compiler:complexity}

Per round, Pass 1 does
$O\!\left(\binom{m}{k} \cdot k \cdot 2^k\right)$ work with the exact
\texttt{optimal2} solver and reduces to $O(m^2)$ at $k{=}2$;
Pass 3 does $O(k \cdot N \cdot q)$ work to scan candidate peer sources
(where $N$ is the number of microstates per lane and $q$ is the buffer
capacity). Pass 4 is linear in the plan size. At Twitter 16p / 2-GPU,
the entire compile takes a few milliseconds and is completely
dominated by I/O overhead elsewhere in the training pipeline.

Every solver uses total orderings on ties (state-index, then lane-id),
so plans are deterministic for a fixed input. This matters for
reproducibility in the evaluation (Section~\ref{sec:eval}) and for the
validator: V1--V5 are strict equalities, so a non-deterministic
compiler would randomly fail otherwise-valid runs.


%-----------------------------------------------------------------------
\section{Runtime Execution}\label{sec:runtime}
% Section 5 -- Runtime Execution

\subsection{From Descriptors to Runtime State}

Phase 3 stopped at descriptor emission: the compiler and validator
produced \texttt{PeerHandoffDescriptor}s, but the runtime still served
every admit from host memory. The final runtime keeps the compiler and
validator unchanged and only changes how those descriptors are
consumed.

At load time, the projected descriptors are handed to
\texttt{MemPartitionBufferStorage}, indexed by
$(\texttt{dst\_lane}, \texttt{round\_idx}, \texttt{partition\_id})$,
and checked for duplicates. Runtime initialization then attempts to
enable CUDA peer access in both directions for every active device
pair. Two environment controls gate the feature:
\texttt{GEGE\_STATEFLOW\_PEER\_RUNTIME=\{off,auto,on\}} controls whether
peer execution is disabled, opportunistic, or required, and
\texttt{GEGE\_STATEFLOW\_PEER\_RUNTIME\_SCOPE=\{all,embeddings\}}
controls whether the peer path is applied to both embeddings and
optimizer state or to embeddings only.

\subsection{In-Place Peer Execution}

Our first Phase 4 prototype allocated a full staged destination buffer
per device. That design was simple but too memory-hungry for
\SI{24}{GB} boards. The final Phase 4b runtime instead performs
\emph{in-place, slot-aligned} peer admits.

At each round boundary, each lane materializes only its own next-state
layout locally, while the runtime consults a precomputed per-round
``globally required'' set to decide which evictions truly leave the
next frontier. This removes the old cross-device next-state barrier.
The source side snapshots only those partitions that will be exported
to other lanes into small per-partition source-scratch tensors and
publishes their slot locations plus a CUDA ready event. Destinations
wait only on the specific source lanes they consume, not on a global
all-lane barrier. This decouples peer reads from subsequent local slot
reuse without allocating another full device-sized buffer.

The destination side then realizes the planner's next-state slot layout
exactly. Retained partitions are permuted into their target slots using
free evict slots as temporary scratch when necessary; partitions that
leave the global next-state frontier are written back to host; and each
new admit is filled either by \texttt{cudaMemcpyPeerAsync} from the
source scratch or by the legacy host path. Because the runtime
materializes the planner's slot layout rather than reusing slots
opportunistically, the compiler's \texttt{src\_slot} and
\texttt{dst\_slot} remain valid at execution time.

\subsection{Safety and Telemetry}

The validator still enforces the structural invariants of
Section~\ref{sec:ir:validator}; the runtime adds dynamic checks that the
projected source lane, destination lane, and slot IDs match the actual
swap-time state. Any mismatch falls back to host reload and increments
a visible counter rather than risking a corrupt swap. This made the
bring-up path auditable: we could distinguish ``descriptor emitted but
not executed,'' ``descriptor executed with fallback,'' and ``descriptor
executed exactly.''

Each epoch reports
\texttt{peer\_bytes\_executed}, \texttt{host\_fallback\_bytes},
\texttt{peer\_copy\_count}, \texttt{descriptor\_mismatch\_count}, and
\texttt{peer\_sync\_wait\_ms}, with per-device breakdowns. These are
the counters used in Section~\ref{sec:eval} to show that the final
Phase 4b runtime executes the planned peer relays with zero fallback
and zero slot mismatch.


%-----------------------------------------------------------------------
\section{Evaluation}\label{sec:eval}
% Section 6 -- Evaluation

\subsection{Experimental Setup}

We evaluate on Twitter (41.6\,M nodes, 1.46\,B edges), 16 partitions,
buffer capacity 4, and two \SI{24}{GB} RTX A5000 GPUs connected over
PCIe. We use the existing GE\textsuperscript{2} Twitter
DistMult/Adagrad configuration and compare three runs:
\textsc{Lane\_Matched} with peer execution disabled,
\textsc{Lane\_Matched} with peer execution enabled, and
\textsc{Disjoint\_Rounds} with peer execution enabled.

For the peer-enabled runs we set
\texttt{GEGE\_STATEFLOW\_PEER\_RUNTIME\_SCOPE=embeddings}, so the
runtime executes peer transfers for node embeddings while optimizer
state remains on the existing host path. Under this scope, executed
peer bytes should equal exactly half of the compiler-reported planned
peer bytes. We report compile-time selection cost, admits, and handoff
distribution from the IR, then runtime epoch time, executed peer
bytes, host fallback bytes, descriptor mismatches, and peer-sync wait
time from the runtime telemetry.

\subsection{Compile-Time Effect}

Table~\ref{tab:ir:example} summarizes the two compiled plans.
\textsc{Lane\_Matched} lowers selection cost from
$1.499 \times 10^{11}$ to $1.291 \times 10^{11}$ (14\%), reduces host
admits from 72 to 62, and changes the handoff distribution from
\{\texttt{FULL\_RELOAD}:10, \texttt{ROTATING\_OVERWRITE}:8,
\texttt{PEER\_RELAY}:8\} to
\{\texttt{FULL\_RELOAD}:0, \texttt{ROTATING\_OVERWRITE}:18,
\texttt{PEER\_RELAY}:18\}. The important point is that this entire
improvement is visible before runtime: the compiler eliminates every
cold admit from the plan and exposes \SI{37.49}{GB} of peer-relay
opportunity, versus \SI{16.66}{GB} for \textsc{Disjoint\_Rounds}.

\subsection{Runtime Effect}

\begin{table}[t]
\centering
\small
\begin{tabular}{lrrrrr}
\toprule
Variant & Epoch (s) & Peer GB & Copies & Fallback GB & Mismatch \\
\midrule
\textsc{Lane\_Matched}, peer off          & 208.256 & 0.00  & 0  & 0.00 & 0 \\
\textsc{Lane\_Matched}, peer on (emb.)    & 209.450 & 18.74 & 18 & 0.00 & 0 \\
\textsc{Disjoint\_Rounds}, peer on (emb.) & 209.525 & 8.33  & 8  & 0.00 & 0 \\
\bottomrule
\end{tabular}
\caption{Twitter 16-partition runtime on two RTX A5000s. Peer-on runs
use \texttt{GEGE\_STATEFLOW\_PEER\_RUNTIME\_SCOPE=embeddings}, so the
executed peer bytes are exactly half of the planned peer bytes in
Table~\ref{tab:ir:example}.}
\label{tab:eval:runtime}
\end{table}

Table~\ref{tab:eval:runtime} gives the runtime result. For
\textsc{Lane\_Matched}, the peer-enabled run executes all 18 planned
embedding handoffs, totaling \SI{18.74}{GB} of device-to-device
traffic, with zero host fallback and zero descriptor mismatch. The
runtime telemetry reports \texttt{peer\_sync\_wait\_ms=1019.646} for
the epoch. For \textsc{Disjoint\_Rounds}, the runtime executes all 8
planned embedding handoffs, totaling \SI{8.33}{GB}, also with zero
fallback and zero mismatch.

\textsc{Lane\_Matched} remains the best peer-enabled schedule, but only
by \SI{0.075}{s}: \SI{209.450}{s} versus \SI{209.525}{s}. Both peer-on
runs are about \SI{1.2}{s} slower than the
\textsc{Lane\_Matched}/peer-off baseline at \SI{208.256}{s}. On this
PCIe A5000 setup, peer execution therefore displaces host traffic
correctly but does not convert that displacement into an epoch-time
speedup. The implication is that the remaining ceiling is the swap
runtime itself rather than descriptor quality: the compiler removes
cold reloads and doubles executed peer traffic, but the wall-clock path
is still dominated by swap orchestration overhead on this hardware.

\subsection{Takeaway}

The runtime result is still valuable. Phase 4b is
\emph{correctness-complete}: the compiler-emitted descriptors survive
projection, validation, slot alignment, and swap-time execution with
zero fallback. The peer-enabled runtime also preserves the compiler's
ordering preference: \textsc{Lane\_Matched} executes more than twice the
peer traffic of \textsc{Disjoint\_Rounds}
(\SI{18.74}{GB} vs.\ \SI{8.33}{GB}) and remains the best peer-enabled
variant. What the result does \emph{not} show is a wall-clock win on
this hardware. The next step is therefore Phase 5 overlap and
staggering, or evaluation on a faster peer interconnect.


%-----------------------------------------------------------------------
\section{Related Work}\label{sec:related}
% Section 7 -- Related Work  [DRAFT STUB]
%
% Clusters:
%   1. Out-of-core KGE / GNN training: GE^2, Marius, DGL-KE, PyTorch
%      BigGraph, BioKG, Edge Partitioning schemes.
%   2. Partition scheduling: RBIBD / balanced incomplete block design
%      literature, bin-packing heuristics for GPU residency.
%   3. Multi-GPU DL training: NCCL all-reduce, ZeRO (DeepSpeed),
%      FSDP -- contrast with out-of-core vs model-parallel regimes.
%   4. IRs for ML: MLIR/Relay/TVM -- borrow "typed IR with cost model"
%      framing; contrast domain (general vs scheduling-only).
%   5. Peer transfers: NCCL P2P, GPUDirect RDMA, prior work on peer
%      memory in DL training (primarily at NVLink-rich servers;
%      less at PCIe-only).

\emph{[Related work pending.]}


%-----------------------------------------------------------------------
\section{Conclusion}\label{sec:conclusion}
Compiled Stateflow gives the multi-GPU out-of-core trainer a
representation for decisions it previously hid: whether an admit is
sourced from host memory, kept in place, or handed across lanes. On
Twitter 16p, the compiler side is already decisive:
\textsc{Lane\_Matched} lowers selection cost by 14\%, cuts host admits
from 72 to 62, and replaces 10 \texttt{FULL\_RELOAD} events with 18
\texttt{PEER\_RELAY} descriptors.

The runtime side is now correct end-to-end. On two \SI{24}{GB} RTX
A5000s, the embeddings-scope peer runtime executes every planned relay
for both \textsc{Lane\_Matched} and \textsc{Disjoint\_Rounds} with zero
host fallback and zero descriptor mismatch, and
\textsc{Lane\_Matched} remains the strongest peer-enabled schedule.
What remains open is performance, not semantics: on this PCIe system,
peer execution displaces host traffic but leaves epoch time effectively
flat relative to the host-reload baseline.

That is still a useful checkpoint. The IR, compiler, validator, and
runtime contract are now stable enough to support the next phase:
bringing peer traffic under optimization with a peer-aware cost model
and reducing swap overhead through more overlapped execution.


%-----------------------------------------------------------------------
\bibliographystyle{ACM-Reference-Format}
\bibliography{references}

\end{document}
